\documentclass[11pt, letterpaper]{article}

\newcommand{\cono}{\tikz[scale=0.15]{
		\draw[fill=orange!30, draw=orange!60!black, line width=0.1pt] (0,0) -- (1,0) -- (0.5,-1.5) -- cycle; 
		\draw[fill=red!50, draw=red!70!black, line width=0.1pt] (0.5,0.25) circle (0.45);
}}

% --- PAQUETES ---
\usepackage[utf8]{inputenc}
\usepackage[T1]{fontenc}
\usepackage[spanish, es-noshorthands]{babel}
\usepackage[top=2cm, bottom=2.2cm, left=2.2cm, right=2.2cm]{geometry}
\usepackage{graphicx}
\usepackage{amsmath}
\usepackage{amssymb}
\usepackage{enumitem}
\usepackage{xcolor}
\usepackage{tikz}
\usetikzlibrary{arrows.meta}
\usepackage{fancyhdr}
\usepackage{eso-pic}
\usepackage[most]{tcolorbox}
\usepackage{booktabs} % Requerido para \toprule, \midrule y \bottomrule

% --- FUENTE SISTEMA ---
\usepackage{helvet}
\renewcommand{\familydefault}{\sfdefault}

% --- COLORES INSTITUCIONALES (LICEO V.A.L.) ---
\definecolor{valblue}{RGB}{0,112,192}
\definecolor{valdark}{RGB}{30,30,30}
\definecolor{valgray}{RGB}{245,247,250}

% Fondo institucional del Liceo V.A.L.
\IfFileExists{../../FondoVAL.pdf}{
	\AddToShipoutPictureBG{
		\AtPageLowerLeft{
			\includegraphics[width=\paperwidth,height=\paperheight,keepaspectratio=false]{../../FondoVAL.pdf}
		}
	}
}{}

% --- CONFIGURACIÓN DE PÁGINA ---
\pagestyle{fancy}
\fancyhf{}
\fancyhead[L]{\footnotesize \color{gray} LICEO V.A.L. -- DEPARTAMENTO DE MATEMÁTICAS}
\fancyhead[R]{\footnotesize \color{gray} MATEMÁTICAS 6º -- TAU 17 (06M-17)}
\fancyfoot[C]{\footnotesize \color{gray} Página \thepage}
\renewcommand{\headrulewidth}{0.4pt}
\renewcommand{\footrulewidth}{0pt}

% --- INTERRUPTOR DINÁMICO DE RESPUESTAS ---
% \showanswerstrue -> Muestra respuestas (Versión Docente)
% \showanswersfalse -> Oculta respuestas (Versión Estudiante)
\newif\ifshowanswers
\showanswerstrue
\newcommand{\answer}[1]{
	\ifshowanswers
	\par\vspace{0.15cm}
	\begin{tcolorbox}[colback=blue!5!white, colframe=blue!60!white, arc=1mm, boxrule=0.5pt, title=\textbf{\footnotesize SOLUCIÓN (DOCENTE)}]
		\color{blue!80!black}\footnotesize #1
	\end{tcolorbox}
	\par\vspace{0.15cm}
	\else
	\par\vspace{0.3cm}
	\fi
}

\begin{document}
	\begin{center}
		
		{\small \textbf{Unidad: TAU 17}}
	\end{center}

	
	
	\begin{enumerate}[leftmargin=*, itemsep=0.8cm]
		
		% --- BLOQUE 1: CONCEPTOS E IDENTIFICACIÓN ---
		\item \textbf{Identificación y Conceptos Fundamentales de Gráficas Estadísticas:}
		\begin{enumerate}[label=\alph*), itemsep=0.25cm]
			\item Observa las opciones del siguiente recuadro y escribe la letra correspondiente (\textbf{A, B, C, D, E} o \textbf{F}) en cada una de las descripciones según el tipo de gráfica que representa:
			\begin{center}
				\begin{tcolorbox}[colback=valblue!10, colframe=valblue, arc=2mm, boxrule=0.8pt, width=0.95\textwidth]
					\centering \small
					\textbf{A.} Gráfica de Barras \quad \textbf{B.} Gráfica Circular (Pastel) \quad \textbf{C.} Pictografía \\
					\textbf{D.} Gráfica de Líneas \quad \textbf{E.} Histograma \quad \textbf{F.} Polígono de Frecuencias
				\end{tcolorbox}
			\end{center}
			\begin{itemize}[itemsep=0.15cm]
				\item [\textbf{[\quad]}] Gráfica que utiliza barras verticales u horizontales separadas entre sí cuya altura representa frecuencias.
				\item [\textbf{[\quad]}] Gráfica circular dividida en sectores que representa porcentajes o partes proporcionalmente de un total.
				\item [\textbf{[\quad]}] Gráfica estadística que utiliza imágenes o íconos representativos para indicar cantidades de datos.
				\item [\textbf{[\quad]}] Gráfica de puntos unidos por segmentos continuos que muestra variaciones o tendencias en el tiempo.
				\item [\textbf{[\quad]}] Gráfica de barras adyacentes (unidas) que representan frecuencias en intervalos numéricos continuos.
				\item [\textbf{[\quad]}] Gráfica formada por segmentos rectos que unen los puntos medios superiores de las barras de un histograma.
			\end{itemize}
			
			
		\end{enumerate}
		
		\answer{
			\textbf{a) Relación de letras:}\\
			- \textbf{A}: Gráfica de barras separadas.\\
			- \textbf{B}: Gráfica circular dividida en sectores.\\
			- \textbf{C}: Pictografía con íconos o dibujos.\\
			- \textbf{D}: Gráfica de líneas continuas.\\
			- \textbf{E}: Histograma con barras unidas por intervalos.\\
			- \textbf{F}: Polígono de frecuencias uniendo puntos medios.\\
			\textbf{b) Análisis teórico:}\\
			- En la gráfica de barras las columnas van separadas (variables cualitativas o discretas), mientras que en el histograma las barras están unidas sin espacios (intervalos continuos).\\
			- Tres árboles y medio: $(3 \times 400) + (0,5 \times 400) = 1.200 + 200 = \mathbf{1.400\text{ kilos}}$.
		}
		
		% --- BLOQUE 2: BARRAS DOBLES ---
		\item \textbf{Lectura de Gráficas de Barras Comparativas -- Cosecha Agrícola:}
		La siguiente gráfica de barras dobles muestra la cosecha de Manzanas y Naranjas (en toneladas) obtenida por una finca durante cuatro trimestres:
		\begin{center}
			\begin{tikzpicture}[scale=0.85]
				\draw[very thin, gray!30] (0,0) grid (9,5);
				\draw[thick, ->] (0,0) -- (9.5,0) node[right] {\footnotesize Trimestre};
				\draw[thick, ->] (0,0) -- (0,5.5) node[above] {\footnotesize Toneladas};
				\foreach \y/\label in {1/50, 2/100, 3/150, 4/200, 5/250} {
					\draw (0,\y) -- (-0.1,\y) node[left] {\footnotesize \label};
				}
				% T1
				\draw[fill=valblue!80, draw=black] (0.8,0) rectangle (1.4, 3.0); % Manzanas: 150
				\draw[fill=orange!70, draw=black] (1.4,0) rectangle (2.0, 2.0); % Naranjas: 100
				\node[below] at (1.4,-0.2) {\footnotesize Trimestre 1};
				% T2
				\draw[fill=valblue!80, draw=black] (2.8,0) rectangle (3.4, 4.4); % Manzanas: 220
				\draw[fill=orange!70, draw=black] (3.4,0) rectangle (4.0, 3.6); % Naranjas: 180
				\node[below] at (3.4,-0.2) {\footnotesize Trimestre 2};
				% T3
				\draw[fill=valblue!80, draw=black] (4.8,0) rectangle (5.4, 2.4); % Manzanas: 120
				\draw[fill=orange!70, draw=black] (5.4,0) rectangle (6.0, 4.0); % Naranjas: 200
				\node[below] at (5.4,-0.2) {\footnotesize Trimestre 3};
				% T4
				\draw[fill=valblue!80, draw=black] (6.8,0) rectangle (7.4, 3.6); % Manzanas: 180
				\draw[fill=orange!70, draw=black] (7.4,0) rectangle (8.0, 2.8); % Naranjas: 140
				\node[below] at (7.4,-0.2) {\footnotesize Trimestre 4};
				% Leyenda
				\fill[valblue!80, draw=black] (1.0, 5.8) rectangle (1.4, 6.1);
				\node[right] at (1.4, 5.95) {\tiny Manzanas};
				\fill[orange!70, draw=black] (4.0, 5.8) rectangle (4.4, 6.1);
				\node[right] at (4.4, 5.95) {\tiny Naranjas};
			\end{tikzpicture}
		\end{center}
		\begin{enumerate}[label=\alph*), itemsep=0.2cm]
			\item ¿En qué trimestre la cosecha de Naranjas superó a la de Manzanas y por cuántas toneladas?
			\item Calcula la producción total acumulada de Manzanas a lo largo de todo el año (suma de los 4 trimestres).
			\item ¿Cuál fue la diferencia entre la producción total de Manzanas y la de Naranjas en el Trimestre 2?
		\end{enumerate}
		
		\answer{
			a) En el \textbf{Trimestre 3} (Naranjas = 200 t, Manzanas = 120 t). La superó por $200 - 120 = \mathbf{80\text{ toneladas}}$.\\
			b) Total Manzanas: $150 + 220 + 120 + 180 = \mathbf{670\text{ toneladas}}$.\\
			c) Diferencia en Trimestre 2: Manzanas (220) - Naranjas (180) = $\mathbf{40\text{ toneladas}}$.
		}
		
		% --- BLOQUE 3: BARRAS HORIZONTALES ---
		\item \textbf{Análisis de Gráficas de Barras Horizontales -- Encuesta Deportiva:}
		Un colegio encuestó a sus estudiantes sobre su deporte preferido, obteniendo los resultados reflejados en la siguiente gráfica de barras horizontales:
		\begin{center}
			\begin{tikzpicture}[yscale=0.65]
				\draw[very thin, gray!30] (0,0) grid (8,5);
				\draw[thick, ->] (0,0) -- (8.5,0) node[right] {\footnotesize Número de Estudiantes};
				\draw[thick, ->] (0,0) -- (0,5.5);
				\foreach \x/\val in {1/20, 2/40, 3/60, 4/80, 5/100, 6/120, 7/140, 8/160} {
					\draw (\x,0) -- (\x,-0.1) node[below] {\footnotesize \val};
				}
				\draw[fill=valblue!70, draw=black] (0,4.2) rectangle (7.5, 4.7);
				\node[right] at (7.6, 4.45) {\footnotesize Fútbol (150)};
				\draw[fill=valblue!70, draw=black] (0,3.2) rectangle (5.0, 3.7);
				\node[right] at (5.1, 3.45) {\footnotesize Baloncesto (100)};
				\draw[fill=valblue!70, draw=black] (0,2.2) rectangle (4.0, 2.7);
				\node[right] at (4.1, 2.45) {\footnotesize Natación (80)};
				\draw[fill=valblue!70, draw=black] (0,1.2) rectangle (2.5, 1.7);
				\node[right] at (2.6, 1.45) {\footnotesize Voleibol (50)};
				\draw[fill=valblue!70, draw=black] (0,0.2) rectangle (1.0, 0.7);
				\node[right] at (1.1, 0.45) {\footnotesize Tenis (20)};
			\end{tikzpicture}
		\end{center}
		\begin{enumerate}[label=\alph*), itemsep=0.2cm]
			\item ¿Cuántos estudiantes en total participaron en la encuesta de deportes?
			\item ¿Cuántos estudiantes más prefieren el Fútbol en comparación con el Baloncesto?
			\item Si se suman los estudiantes que eligieron Voleibol y Tenis, ¿alcanzan la cantidad de los que eligieron Natación? Justifica con cifras.
		\end{enumerate}
		
		\answer{
			a) Total de estudiantes encuestados: $150 + 100 + 80 + 50 + 20 = \mathbf{400\text{ estudiantes}}$.\\
			b) Diferencia entre Fútbol y Baloncesto: $150 - 100 = \mathbf{50\text{ estudiantes más}}$.\\
			c) Voleibol + Tenis = $50 + 20 = 70$ estudiantes. No alcanzan la cifra de Natación ($80$ estudiantes), ya que $70 < 80$.
		}
		
		%\newpage
		% --- BLOQUE 4: LÍNEAS Y MOVIMIENTO ---
		\item \textbf{Gráfica de Líneas Continua -- Competencia en Bicicleta:}
		La siguiente gráfica de líneas representa el recorrido en kilómetros de dos ciclistas (Carlos y Sofía) durante una prueba de entrenamiento a lo largo del tiempo:
		\begin{center}
			\begin{tikzpicture}[scale=0.95]
				\draw[very thin, gray!30] (0,0) grid (7,5);
				\draw[thick, ->] (0,0) -- (7.5,0) node[right] {\footnotesize Tiempo (Horas)};
				\draw[thick, ->] (0,0) -- (0,5.5) node[above] {\footnotesize Distancia (km)};
				\foreach \x/\label in {1/1h, 2/2h, 3/3h, 4/4h, 5/5h, 6/6h} {
					\draw (\x,0) -- (\x,-0.1) node[below] {\footnotesize \label};
				}
				\foreach \y/\val in {1/10, 2/20, 3/30, 4/40, 5/50} {
					\draw (0,\y) -- (-0.1,\y) node[left] {\footnotesize \val};
				}
				% Carlos (Rojo continuo)
				\draw[ultra thick, red!80!black] (0,0) -- (2, 3.0) -- (3.5, 3.0) -- (6, 5.0);
				% Sofía (Azul punteado)
				\draw[ultra thick, valblue, dashed] (0,0) -- (2, 1.5) -- (3.5, 3.0) -- (6, 5.0);
				\node[red!80!black, right] at (6, 4.7) {\small Carlos};
				\node[valblue, right] at (6, 5.3) {\small Sofía};
			\end{tikzpicture}
		\end{center}
		\begin{enumerate}[label=\alph*), itemsep=0.2cm]
			\item ¿Cuántos kilómetros había recorrido Carlos en las primeras 2 horas y qué ocurrió con él entre las 2h y las 3,5h?
			\item ¿En qué hora exacta Sofía alcanzó a Carlos y cuántos kilómetros habían recorrido ambos en ese punto?
			\item ¿Cuál de los dos ciclistas ganó la prueba al llegar primero o cómo terminaron al cabo de 6 horas?
		\end{enumerate}
		
		\answer{
			a) En las primeras 2 horas Carlos recorrió \textbf{$30\text{ km}$}. Entre las 2h y 3,5h la gráfica es horizontal, lo que significa que \textbf{se detuvo a descansar} (la distancia se mantuvo constante en $30\text{ km}$).\\
			b) Sofía alcanzó a Carlos exactamente a las \textbf{3,5 horas}, cuando ambos estaban a una distancia de \textbf{$30\text{ km}$}.\\
			c) Ambos recorrieron $50\text{ km}$ en 6 horas y llegaron al mismo tiempo, de modo que \textbf{empataron al final del recorrido}.
		}
		

		
		% --- BLOQUE 6: PORCENTAJES Y BARRAS ---
		\item \textbf{Barras Horizontales y Porcentajes -- Reciclaje de Residuos:}
		Se midió el porcentaje de residuos sólidos reciclados en cinco municipios durante el último año:
		\begin{center}
			\begin{tikzpicture}[yscale=0.6]
				\draw[very thin, gray!30] (0,0) grid (6,5);
				\draw[thick, ->] (0,0) -- (6.5,0) node[right] {\footnotesize \% de Residuos Reciclados};
				\draw[thick, ->] (0,0) -- (0,5.5);
				\foreach \x/\val in {1/10\%, 2/20\%, 3/30\%, 4/40\%, 5/50\%} {
					\draw (\x,0) -- (\x,-0.1) node[below] {\footnotesize \val};
				}
				\draw[fill=valblue!70, draw=black] (0,4.2) rectangle (4.5, 4.7);
				\node[left] at (-0.1, 4.45) {\small Municipio A};
				\node[right] at (4.5, 4.45) {\small 45\%};
				\draw[fill=valblue!70, draw=black] (0,3.2) rectangle (3.0, 3.7);
				\node[left] at (-0.1, 3.45) {\small Municipio B};
				\node[right] at (3.0, 3.45) {\small 30\%};
				\draw[fill=valblue!70, draw=black] (0,2.2) rectangle (1.8, 2.7);
				\node[left] at (-0.1, 2.45) {\small Municipio C};
				\node[right] at (1.8, 2.45) {\small 18\%};
				\draw[fill=valblue!70, draw=black] (0,1.2) rectangle (2.5, 1.7);
				\node[left] at (-0.1, 1.45) {\small Municipio D};
				\node[right] at (2.5, 1.45) {\small 25\%};
				\draw[fill=valblue!70, draw=black] (0,0.2) rectangle (1.0, 0.7);
				\node[left] at (-0.1, 0.45) {\small Municipio E};
				\node[right] at (1.0, 0.45) {\small 10\%};
			\end{tikzpicture}
		\end{center}
		\begin{enumerate}[label=\alph*), itemsep=0.2cm]
			\item ¿Cuál es el municipio con el mejor porcentaje de reciclaje y cuál es el de menor porcentaje?
			\item ¿En cuántos puntos porcentuales supera el Municipio A al Municipio D?
		\end{enumerate}
		
		\answer{
			a) El mejor es el \textbf{Municipio A} (45\%) y el de menor porcentaje es el \textbf{Municipio E} (10\%).\\
			b) Diferencia entre A y D: $45\% - 25\% = \mathbf{20\text{ puntos porcentuales}}$.
		}
		
		%\newpage
		% --- BLOQUE 7: LÍNEAS DOBLES / TENDENCIA HISTÓRICA ---
		\item \textbf{Líneas Dobles -- Comparación de Ventas Tecnológicas:}
		La gráfica muestra el historial de ventas anuales (en millones de dólares) de dos tiendas de tecnología (Tienda Alfa y Tienda Beta) entre 2018 y 2023:
		\begin{center}
			\begin{tikzpicture}[scale=0.9]
				\draw[very thin, gray!30] (0,0) grid (6,5);
				\draw[thick, ->] (0,0) -- (6.5,0) node[right] {\footnotesize Año};
				\draw[thick, ->] (0,0) -- (0,5.5) node[above] {\footnotesize Ventas (Millones \$)};
				\foreach \x/\label in {1/2018, 2/2019, 3/2020, 4/2021, 5/2022, 6/2023} {
					\draw (\x,0) -- (\x,-0.1) node[below] {\footnotesize \label};
				}
				\foreach \y/\val in {1/2, 2/4, 3/6, 4/8, 5/10} {
					\draw (0,\y) -- (-0.1,\y) node[left] {\footnotesize \val};
				}
				% Alfa (Azul)
				\draw[ultra thick, valblue, mark=*] plot coordinates { (1, 1.5) (2, 2.5) (3, 2.0) (4, 3.5) (5, 4.2) (6, 5.0) };
				% Beta (Rojo)
				\draw[ultra thick, red!80!black, mark=square*] plot coordinates { (1, 1.0) (2, 1.8) (3, 3.0) (4, 3.0) (5, 3.8) (6, 4.1) };
				\node[valblue, right] at (6, 5.0) {\small Alfa (\$10M)};
				\node[red!80!black, right] at (6, 4.1) {\small Beta (\$8.2M)};
			\end{tikzpicture}
		\end{center}
		\begin{enumerate}[label=\alph*), itemsep=0.2cm]
			\item ¿En qué único año la Tienda Beta tuvo mayores ventas que la Tienda Alfa?
			\item ¿De cuánto fueron las ventas totales de la Tienda Alfa en el año 2023?
			\item ¿Entre 2020 y 2023, qué tienda tuvo un crecimiento constante sin caídas?
		\end{enumerate}
		
		\answer{
			a) En el año \textbf{2020}, donde la Tienda Beta vendió $\$6\text{M}$ y Alfa cayó a $\$4\text{M}$.\\
			b) En 2023 la Tienda Alfa vendió \textbf{$\$10\text{ millones de dólares}$}.\\
			c) Ambas tiendas crecieron entre 2020 y 2023, pero la \textbf{Tienda Alfa} tuvo un crecimiento sostenido pasando de $\$4\text{M}$ a $\$10\text{M}$.
		}
		\newpage 
		\item Una persona desea enviar un paquete de 3 lb a la Zona 4 de correo prioritario, y otro paquete de 1 lb a la Zona 3. Utilizando la siguiente tabla de tarifas postales de 2017 de tu libro de texto, calcula el costo total del envío de ambos paquetes:
		
		\begin{center}
			\begin{tabular}{ccccc}
				\toprule
				\textbf{Peso} & \textbf{Zonas 1 y 2} & \textbf{Zona 3} & \textbf{Zona 4} & \textbf{Zona 5} \\
				\midrule
				1 lb & \$5,75 & \$6,26 & \$6,37 & \$6,53 \\
				2 lb & \$6,33 & \$6,40 & \$6,63 & \$7,91 \\
				3 lb & \$6,41 & \$7,16 & \$7,62 & \$9,22 \\
				4 lb & \$6,62 & \$7,39 & \$8,46 & \$10,19 \\
				\bottomrule
			\end{tabular}
		\end{center}
		
		\answer{
			Buscamos los valores en la tabla por intersección de fila (peso) y columna (zona):\\
			- Paquete de 3 lb enviado a la Zona 4: \textbf{\$7,62}\\
			- Paquete de 1 lb enviado a la Zona 3: \textbf{\$6,26}\\
			Costo total = \$7,62 + \$6,26 = \textbf{\$13,88}
		}
		
		
		\item En la heladería \textit{Barney's Café}, se registró la venta diaria de helados de diferentes poblaciones a través de un pictograma. Si cada símbolo \cono\ equivale a \textbf{\$100}, ¿cuánto dinero en total gastaron los niños y los ancianos juntos en helados?
		
		\begin{center}
			\begin{tabular}{rl}
				\toprule
				\textbf{Grupo de Clientes} & \textbf{Ventas de Helados (Cada \cono\ = \$100)} \\
				\midrule
				Niños: & \cono\ \cono\ \cono\ \cono\ \cono \\
				Padres: & \cono\ \cono\ \cono\ \cono \\
				Ancianos: & \cono\ \cono \\
				\bottomrule
			\end{tabular}
		\end{center}
		
		\answer{
			- Contamos los símbolos para cada uno de los grupos indicados:\\
			* Niños: 5 símbolos $\rightarrow$ 5 $\times$ \$100 = \$500.\\
			* Ancianos: 2 símbolos $\rightarrow$ 2 $\times$ \$100 = \$200.\\
			- Sumamos el total de dinero gastado por ambos grupos: \$500 + \$200 = \textbf{\$700}.\\
			*(O de forma equivalente: 5 + 2 = 7 símbolos en total; 7 $\times$ \$100 = \textbf{\$700}).
		}
	\end{enumerate}
\end{document}